\documentclass[11pt,a4paper]{article}
\usepackage[margin=2.5cm]{geometry}
\usepackage{amsmath,amssymb,amsthm}
\usepackage{booktabs}
\usepackage{graphicx}
\usepackage{caption}
\usepackage[colorlinks=true,linkcolor=blue,citecolor=blue,urlcolor=blue]{hyperref}
\usepackage[round,authoryear]{natbib}

\theoremstyle{plain}
\newtheorem{proposition}{Proposition}
\newtheorem{lemma}{Lemma}
\newtheorem{corollary}{Corollary}

\newcommand{\TI}{\mathrm{TI}}
\newcommand{\CV}{\mathrm{CV}}
\newcommand{\rms}{\mathrm{rms}}
\newcommand{\sd}{\mathrm{sd}}

\title{\bfseries Degree-based topological indices measure molecular size:\\
a null-baseline protocol for QSPR}
\author{Eric Schultz\\[2pt]
\small Independent Researcher \quad ORCID 0009-0006-6283-1696}
\date{Draft --- prepared for submission to \emph{Journal of Cheminformatics}}

\begin{document}
\maketitle

\begin{abstract}
\noindent
In QSPR studies using degree-based topological indices, the reported coefficient of
determination and the descriptor's actual contribution move in \emph{opposite}
directions: as $R^2$ rises from $0.34$ to $0.94$, the improvement over simply counting
heavy atoms collapses from $+0.148$ to $+0.005$. We explain this. Writing a degree-based
index as $\TI_f(G)=m\langle p(G),f\rangle$, with $m$ the edge count and $p(G)$ the
edge-degree-pair profile, we bound the correlation between any such index and $m$ in
terms of three quantities: the variation of the weighting function along directions in
which real molecular profiles differ, the mean edge weight, and the coefficient of
variation of molecular size in the sample. A corollary quantified over \emph{all}
properties follows: no degree-based index can differ from edge count in its correlation
with \emph{any} property by more than $\sqrt{2(1-\rho)}$, where $\rho$ is its correlation
with edge count. On $8{,}583$ drug-like molecules the bound is valid for all 14 indices
tested and tight to within $0.013$ for ten; the corollary holds without exception across
70 index--property pairs. Empirically the 14 indices span one dimension
(PC1 $=97.5\%$ on 6{,}000 ZINC molecules). Against molecular weight no index improves on
heavy-atom count on any dataset tested; against measured melting points the best single
index reaches $R^2=0.157$, below a two-number baseline of atom count plus ring count
($0.218$). At $N\approx15$ a permutation null places the 95th percentile at $R^2\approx0.35$
for linear fits, $0.50$ for quadratic and $0.63$ for cubic --- so raising the fit order
raises the bar a claim must clear rather than lowering it. We recommend a referee-facing
checklist whose central element is a single diagnostic: report $\Delta R^2$, the
improvement over a heavy-atom-count null.
\end{abstract}

\noindent\textbf{Keywords:} topological index; QSPR; molecular descriptor; degree-based
index; null baseline; chemical graph theory

\section{Introduction}

A molecular graph $G$ represents a hydrogen-suppressed molecule with atoms as vertices and
bonds as edges. A \emph{degree-based topological index} assigns a real number to $G$ by
summing a symmetric function of the endpoint degrees over all edges,
\[
\TI_f(G)=\sum_{uv\in E(G)} f(d_u,d_v).
\]
Instances include the Zagreb indices \citep{gutman1972}, the Randi\'c connectivity index
\citep{randic1975}, the atom--bond connectivity (ABC) and geometric--arithmetic (GA)
indices, the forgotten index \citep{furtula2015}, and the Sombor index
\citep{gutman2021}. Several thousand such invariants have been proposed.

\subsection{The template under examination}

A recognisable publication template has become common in pharmacy,
pharmaceutical-chemistry and applied-mathematics journals: select a therapeutic class and
10--25 drug molecules; compute 6--20 degree-based indices; obtain 5--13 physicochemical
properties, typically from an online compilation; regress each index against each
property; and report, per property, the index with the largest $R^2$, described as the
best predictor. The properties most often used are molecular weight, molar refractivity,
molar volume, polarizability, surface tension, boiling point and flash point.

\subsection{The difficulty}

Four of those properties --- molecular weight, molar refractivity, molar volume,
polarizability --- are \emph{additive}: each is, to good approximation, a sum of atomic or
group contributions, and hence a near-linear function of molecular size. Degree-based
indices are also near-linear functions of size. A high $R^2$ between the two is therefore
close to guaranteed and carries little information about the index.

This is not a new suspicion. \citet{gutman2013} tested 20 vertex-degree-based indices
against heats of formation and normal boiling points of \emph{octane isomers} and found
the correlation ability of many to be weak or nil, with augmented Zagreb and ABC
performing best. Their design fixes $n$ and $m$, removing size variation entirely --- and,
we note, forcing the choice of non-additive properties, since molecular weight and molar
refractivity are literally constant within a constitutional-isomer class. That paper is
cited in a large fraction of the studies described above, often as introductory
boilerplate, without its implication being drawn.

\citet{randic1994} had earlier addressed descriptor intercorrelation directly, arguing
that apparent parallelism need not be a concern unless descriptors are strictly collinear,
and that orthogonalization reveals genuine structural differences even between nearly
collinear descriptors. We accept this as stated. We observe only that it concerns
\emph{orthogonalized combinations}. It provides no licence for the comparative claim
--- ``index $X$ is the best predictor of property $Y$'' --- which the template produces.

\section{Theory}

\subsection{Profile decomposition}

Let $P=\{(i,j):1\le i\le j\le \Delta\}$ be the set of unordered degree pairs available at
maximum degree $\Delta$ (for hydrogen-suppressed organic molecules $\Delta=4$, so
$|P|=10$). For a molecular graph $G$ with $m$ edges define the \emph{edge-degree-pair
profile} $p(G)\in\Delta(P)$ by
\[
p_\pi(G)=\frac{\#\{uv\in E(G):\{d_u,d_v\}=\pi\}}{m},
\]
a point in the probability simplex over $P$. Regarding $f$ as a vector in $\mathbb{R}^P$,
every degree-based index factorises exactly:
\begin{equation}
\TI_f(G)=m(G)\cdot\langle p(G),f\rangle .
\label{eq:factor}
\end{equation}
The index is the product of an \emph{extensive} factor (edge count) and an
\emph{intensive} factor (mean edge weight). All structural information beyond size is
carried by $\langle p(G),f\rangle$.

\subsection{The bound}

Consider molecules $G_1,\dots,G_N$ with edge counts $m_i$ and profiles $p_i$. Let
$\bar p=N^{-1}\sum_i p_i$, $u_i=p_i-\bar p$, $\bar g=\langle\bar p,f\rangle>0$ and
$\CV=\sd(m)/\mathrm{mean}(m)$. Write $\tilde v$ for the mean-centred version of
$v\in\mathbb{R}^N$.

\begin{proposition}
Let $E_i=m_i\langle u_i,f\rangle$ and $\tau_f=\|\tilde E\|/(\bar g\|\tilde m\|)$. If
$\tau_f\le1$ then
\begin{equation}
r^2(\TI_f,m)\;\ge\;1-\tau_f^2 .
\label{eq:bound}
\end{equation}
Moreover, with $\Sigma=\sum_i w_i u_iu_i^{\!\top}$, $w_i=m_i^2/\sum_j m_j^2$, and
$f_c=f-c\mathbf{1}$ for any scalar $c$,
\begin{equation}
\tau_f\;\le\;\frac{\sqrt{f_c^{\!\top}\Sigma f_c}\;\rms(m)}{\bar g\,\sd(m)} .
\label{eq:spectral}
\end{equation}
\end{proposition}

\begin{proof}
From \eqref{eq:factor}, $\TI_f(G_i)=\bar g m_i+E_i$. Set $S=\bar g m$, so
$\TI_f=S+E$ and, centring being linear, $\widetilde{\TI_f}=\tilde S+\tilde E$ with
$\tilde S=\bar g\tilde m$ parallel to $\tilde m$.

For \eqref{eq:bound}, note $r(\TI_f,m)=\cos\angle(\widetilde{\TI_f},\tilde m)
=\cos\angle(\widetilde{\TI_f},\tilde S)$. The set $\{\tilde S+e:\|e\|\le\rho\}$ is a
Euclidean ball of radius $\rho$ centred at $\tilde S$. Every point of that ball lies within
angle $\arcsin(\rho/\|\tilde S\|)$ of $\tilde S$ when $\rho<\|\tilde S\|$, the extreme
being attained on the tangent line from the origin. Taking $\rho=\|\tilde E\|$ gives
$\angle(\widetilde{\TI_f},\tilde S)\le\arcsin\tau_f$ and hence
$r\ge\sqrt{1-\tau_f^2}$.

For \eqref{eq:spectral}: each $u_i$ satisfies $\langle u_i,\mathbf{1}\rangle=0$, so
$\langle u_i,f\rangle=\langle u_i,f_c\rangle$ for every $c$ --- the weighting function is
determined only up to an additive constant. Centring is an orthogonal projection, so
$\|\tilde E\|\le\|E\|$, and
\[
\|E\|^2=\sum_i m_i^2\langle u_i,f_c\rangle^2
=\Big(\sum_j m_j^2\Big)\,f_c^{\!\top}\Sigma f_c ,
\]
whence $\|E\|=\sqrt{N}\,\rms(m)\sqrt{f_c^{\!\top}\Sigma f_c}$. Dividing by
$\bar g\|\tilde m\|=\bar g\sqrt{N}\,\sd(m)$ gives \eqref{eq:spectral}.
\end{proof}

\subsection{A corollary over all properties}

The Proposition constrains the relationship between $\TI_f$ and edge count; the claim of
interest concerns properties. The gap closes in one step.

\begin{lemma}
For any random variables $X,Y,Z$,
$\;|r(X,Z)-r(Y,Z)|\le\sqrt{2\bigl(1-r(X,Y)\bigr)}$.
\end{lemma}

\begin{proof}
Let $x,y,z$ be the centred, unit-normalised data vectors. Then
$r(X,Z)-r(Y,Z)=\langle x-y,z\rangle\le\|x-y\|\,\|z\|=\|x-y\|$, and
$\|x-y\|^2=2-2\langle x,y\rangle=2(1-r(X,Y))$. Symmetry gives the absolute value.
\end{proof}

\begin{corollary}
Let $\rho=r(\TI_f,m)\ge\sqrt{1-\tau_f^2}$ by \eqref{eq:bound}. Then for \emph{every}
property $P$ whatsoever,
\begin{equation}
\bigl|r(\TI_f,P)-r(m,P)\bigr|\;\le\;\sqrt{2(1-\rho)}\;\le\;
\sqrt{2\bigl(1-\sqrt{1-\tau_f^2}\bigr)} .
\label{eq:corollary}
\end{equation}
\end{corollary}

This is the operative statement. It is not a claim about the properties we happened to
test: for any response variable at all --- measured, computed, physicochemical,
biological --- a degree-based index and a bond count are interchangeable to within a
margin fixed by $\tau_f$, a property of the index and the molecule set alone. For ABC on
our melting-point set, $\rho=0.9983$ and the margin is $0.058$; across 14 indices and 5
properties there are no violations and the largest observed deviation for ABC is $0.015$.

\subsection{Interpretation}

The bound factorises into three parts, each with a chemical reading.
\begin{itemize}\itemsep2pt
\item $\sqrt{f_c^{\!\top}\Sigma f_c}$ --- \textbf{index spread along realised variation}:
not how much $f$ varies over all conceivable degree pairs, but how much it varies along
the directions in which real molecular profiles actually differ. This is why the naive
sandwich $f_{\min}m\le\TI_f\le f_{\max}m$, which uses the worst-case range of $f$, is far
too weak: it ignores the covariance structure entirely.
\item $1/\bar g$ --- \textbf{normalisation}. Only the relative spread of $f$ matters.
\item $\rms(m)/\sd(m)\approx\sqrt{1+\CV^2}/\CV$ --- \textbf{size dispersion}, which
diverges as $\CV\to0$.
\end{itemize}
These are the complete list of things one could break to construct a degree-based index
that is not a size proxy. Two are properties of the index; the third is a property of the
\emph{study design}, and is under the analyst's control.

The $\CV$ term recovers \citet{gutman2013} immediately. On an isomer set $\CV=0$ exactly,
the bound is vacuous, and the indices are free to differ arbitrarily from edge count and
from one another --- precisely the regime in which comparing them is meaningful. Their
study and the studies of \S1.1 were never in conflict; they sit at opposite ends of a
continuum that \eqref{eq:bound}--\eqref{eq:spectral} parameterises.

\section{Methods}

\paragraph{Index battery.} Fourteen degree-based indices with weighting functions
$f(a,b)$: first Zagreb $a+b$; second Zagreb $ab$; forgotten $a^2+b^2$; hyper-Zagreb
$(a+b)^2$; Randi\'c $(ab)^{-1/2}$; sum-connectivity $(a+b)^{-1/2}$; ABC
$\sqrt{(a+b-2)/(ab)}$; geometric--arithmetic $2\sqrt{ab}/(a+b)$; harmonic $2/(a+b)$;
Sombor $\sqrt{a^2+b^2}$; augmented Zagreb $\bigl(ab/(a+b-2)\bigr)^3$; inverse sum indeg
$ab/(a+b)$; symmetric division deg $a/b+b/a$; redefined third Zagreb $ab(a+b)$.

\paragraph{Datasets.} A hand-curated drug set ($N=38$) from the therapeutic classes
recurring in the target literature; a random subsample of a drug-like ZINC collection
($N=6{,}000$); ESOL measured aqueous solubility ($N=1{,}105$, \citealp{delaney2004});
Open Notebook Science melting points ($N=8{,}583$); and a boiling-point compilation
($N=6{,}400$). Structures for the last two were joined by CAS lookup through PubChem.
Filters: 10--60 heavy atoms, maximum degree $\le 4$, must contain carbon, salts and
mixtures excluded.

\paragraph{Null baseline.} The null descriptor is the \textbf{heavy-atom count} $n$. For
each property we report $R^2(n)$, $R^2$ of the best of the 14 indices, and
$\Delta R^2=R^2(\text{best})-R^2(n)$. Out-of-sample performance uses repeated 70/30
splits. Small-sample behaviour uses repeated random draws of 15 molecules, matching the
modal published study size.

\section{Results}

\subsection{The indices span one dimension}

\begin{table}[h]\centering\small
\caption{Dimensionality of the 14-index battery across two chemical spaces.}
\begin{tabular}{lrr}
\toprule
& drug set ($N=38$) & ZINC ($N=6{,}000$)\\
\midrule
PC1 variance explained & 98.64\% & 97.49\%\\
PC2 & 1.13\% & 1.76\%\\
PC3 & 0.21\% & 0.69\%\\
median pairwise $|r|$ & 0.989 & 0.979\\
minimum pairwise $|r|$ & 0.934 & 0.908\\
$\min r(\TI,m)$ & 0.961 & 0.947\\
\bottomrule
\end{tabular}
\end{table}

\begin{figure}[h]\centering
\includegraphics[width=\textwidth]{figure3_dimensionality.png}
\caption{(a) Pairwise correlations among the 14 indices on 8{,}583 drug-like molecules.
(b) Scree plot on a logarithmic axis: the battery has one effective dimension.}
\end{figure}

The result replicates across unrelated chemical spaces and is not an artefact of size
range: the ZINC sample has roughly half the size dispersion of the drug set
($\CV(m)=0.213$ vs $0.408$) yet shows the same collapse. A rank statement makes the point
exactly. On ESOL, cross-validated $R^2$ for measured $\log S$: $n$ alone $0.3199$; best
single index $0.3348$; all 14 indices $0.3803$; $n$ plus all 14 indices $0.3803$ ---
identical to four decimal places. Heavy-atom count already lies in their span.

The observed edge-degree-pair profile is highly concentrated: across the drug set four of
the ten available pairs account for 91\% of all edges, namely $(2,3)$ 41.3\%, $(2,2)$
17.7\%, $(3,3)$ 17.5\% and $(1,3)$ 14.2\%.

\subsection{The bound is valid and tight}

\begin{figure}[h]\centering
\includegraphics[width=0.62\textwidth]{figure4_bound.png}
\caption{Spectral bound \eqref{eq:spectral} against observed $r^2(\TI_f,m)$ on the
melting-point set. All 14 points lie above the diagonal (bound valid); ten lie within
$0.013$ of it. Per-index values are given in the Appendix.}
\end{figure}

The spectral bound is valid in 14/14 cases and non-vacuous in 14/14. A H\"older bound
using $\|u_i\|_1$ and the range of $f$ is valid but vacuous ($\tau>1$) for 6 of the 14
indices: the covariance structure is essential.

\subsection{Against additive properties, no index beats counting atoms}

We anchor on \textbf{molecular weight}, which is exactly additive and is not the output of
any fitted model. Molar refractivity (Crippen) and Labute ASA are themselves fitted
additive schemes, so correlations with them are partly circular; we report them for
comparability with the literature but they carry no independent weight.

\begin{table}[h]\centering\small
\caption{Null-baseline comparison on additive properties. Negative $\Delta R^2$ means the
heavy-atom count beats every one of the 14 indices.}
\begin{tabular}{llrrr}
\toprule
dataset & property & $R^2(n)$ & $R^2$(best) & $\Delta R^2$\\
\midrule
Drug set (38) & molecular weight & 0.9850 & 0.9826 & $-0.0025$\\
              & Labute ASA & 0.9911 & 0.9930 & $+0.0019$\\
              & molar refractivity & 0.9564 & 0.9696 & $+0.0131$\\
ESOL (1{,}105) & molecular weight & 0.8510 & 0.8645 & $+0.0136$\\
              & molar refractivity & 0.9414 & 0.9432 & $+0.0018$\\
              & Labute ASA & 0.9618 & 0.9521 & $-0.0096$\\
$T_m$ set (8{,}583) & molecular weight & 0.7977 & 0.7941 & $-0.0036$\\
              & molar refractivity & 0.8856 & 0.9128 & $+0.0272$\\
$T_b$ set (6{,}400) & molecular weight & 0.8430 & 0.8136 & $-0.0295$\\
              & molar refractivity & 0.8233 & 0.8824 & $+0.0591$\\
\bottomrule
\end{tabular}
\end{table}

In four of ten cases the null beats every index outright; the largest gain anywhere is
$+0.059$.

\subsection{Against measured non-additive properties}

No single index exceeds $R^2=0.212$ on a measured non-additive property. We state the
countervailing result plainly: used \emph{jointly}, the 14 indices do better than any of
them alone.

\begin{table}[h]\centering\small
\caption{Cross-validated $R^2$ for measured melting point ($N=8{,}583$).}
\begin{tabular}{lr}
\toprule
model & CV $R^2$\\
\midrule
heavy-atom count $n$ & 0.0713\\
molecular weight & 0.0626\\
best single index (ReZG3) & 0.1545\\
\textbf{$n$ + ring count} & \textbf{0.2179}\\
all 14 indices & 0.3364\\
all 14 indices $+\,n\,+$ ring count & 0.3359\\
\bottomrule
\end{tabular}
\end{table}

Two things follow. The best single index is beaten by a two-number baseline requiring no
chemical graph theory at all. And the 14 indices jointly \emph{do} carry information
beyond simple counts ($0.336$ vs $0.218$) --- a real effect we do not dispute --- but
adding $n$ and ring count to the battery changes cross-validated performance by
$-0.0005$, confirming that trivial size descriptors already lie within their span. The
joint model is a legitimate descriptor set; the \emph{comparative single-index ranking}
that the template produces is not a legitimate inference.

\subsection{Reported $R^2$ is a readout of sample size dispersion}

\begin{figure}[h]\centering
\includegraphics[width=\textwidth]{size_dispersion.png}
\caption{Reported $R^2$ as a function of the size dispersion of the sample, for two
additive properties and one measured property. Subsamples of 15 molecules from ESOL. Red:
best of 14 indices; black dashed: heavy-atom-count null. Shaded band: interquartile range.
Orange region marks the dispersion typical of published drug sets.}
\end{figure}

\begin{figure}[h]\centering
\includegraphics[width=\textwidth]{figure2.png}
\caption{(a) The inversion: as the headline $R^2$ rises with size dispersion, the
improvement over the atom-count null falls toward zero. (b) Sampling distribution of the
reported statistic (filled) against its true permutation null (outline).}
\end{figure}

Nothing changes across a row of Figure~3 except how much size variation the sample
contains. As the headline $R^2$ for molar refractivity climbs from $0.34$ to $0.94$, the
index's contribution over atom count collapses from $+0.148$ to $+0.005$: \emph{the most
impressive-looking results are those in which the descriptor is doing the least work.}

At published sample sizes the winner is unstable --- all 14 indices win at some frequency,
the modal winner taking 42--48\% for melting and boiling point and as little as 14\% at
low size dispersion. Contrasting the sampling distribution with the permutation null
(Figure~4b) gives two conclusions we are careful to separate. \textbf{(i)} Any best-of-14
$R^2$ below roughly $0.35$ at $N=15$ falls inside the permutation null and is not
evidence of a structure--property relationship; a substantial share of published values
lies in this range. \textbf{(ii)} Values well above that are \emph{not} explicable as
noise --- they are explicable by upward bias. At $N=15$ the expected best-of-14 $R^2$ is
$0.306$ against a full-sample value of $0.212$, with a 5th--95th range of $0.026$--$0.730$.
A single reported value from a 15-molecule study is compatible with almost any underlying
effect size, and the publication filter selects the upper tail. We do not claim published
results are noise; we claim they are uninformative as point estimates.

\subsection{Functional form: the threshold is not fixed}

Studies in the target literature frequently fit quadratic or cubic rather than linear
models. We repeated the analysis at polynomial degrees 1--3, fitting the \emph{same} order
to the index and to the null.

At full sample size the verdict is unchanged, with one honest exception. For molecular
weight $\Delta R^2$ is $-0.0036$, $-0.0036$ and $-0.0041$ at degrees 1, 2 and 3: the null
beats every index at every order. For boiling point $\Delta R^2$ is likewise flat
($+0.060$, $+0.057$, $+0.057$). For melting point, however, the indices benefit from
curvature more than atom count does and $\Delta R^2$ \emph{rises} from $+0.084$ to
$+0.112$; our linear-only analysis was mildly conservative there. In absolute terms both
remain poor ($R^2=0.085$ for the null against $0.197$ for the best index at degree 3), and
cross-validated values track in-sample values to within $0.001$ at $N=8{,}583$.

At $N=15$ the effect is dramatic. Fitting a cubic means estimating four parameters from
fifteen points, and the permutation null inflates accordingly.

\begin{table}[h]\centering\small
\caption{Permutation null and real-draw distributions of best-of-14 $R^2$ at $N=15$.}
\begin{tabular}{llrrrr}
\toprule
property & fit & null mean & \textbf{null 95th} & real mean & real 95th\\
\midrule
melting point & linear & 0.104 & \textbf{0.323} & 0.248 & 0.582\\
              & quadratic & 0.207 & \textbf{0.475} & 0.351 & 0.664\\
              & cubic & 0.316 & \textbf{0.605} & 0.442 & 0.729\\
boiling point & linear & 0.117 & \textbf{0.358} & 0.311 & 0.734\\
              & quadratic & 0.237 & \textbf{0.533} & 0.435 & 0.782\\
              & cubic & 0.349 & \textbf{0.654} & 0.543 & 0.848\\
\bottomrule
\end{tabular}
\end{table}

\begin{figure}[h]\centering
\includegraphics[width=0.68\textwidth]{figure5_fitorder.png}
\caption{Raising the polynomial order of the fitted model raises the evidential bar. A
claim must clear the red line to be distinguishable from a permutation null at the same
sample size and fit order.}
\end{figure}

The evidential threshold is therefore \emph{not a fixed number but a function of the fit
order the study used}. A paper reporting $R^2=0.60$ from a cubic fit to 15 molecules has
produced a value lying inside its own permutation null; the same value from a linear fit
would be well outside it. Higher-order fits do not strengthen a small-sample QSPR claim
--- they raise the bar it must clear, and papers reporting improved $R^2$ after moving
from linear to cubic at $N\approx15$ are reporting overfitting.

\section{Discussion}

\subsection{What is and is not claimed}

We do \textbf{not} claim degree-based indices carry no information. A selection-bias null
(shuffled response, maximum over 14 indices) gives mean $R^2=0.0011$ and 95th percentile
$0.0040$ on ESOL, against an observed $\Delta R^2$ of $+0.0143$: a small but real margin.
Out of sample the indices genuinely help on some additive properties.

What we claim is narrower: \textbf{the comparative claims produced by this template are
not interpretable}, because the objects being compared span one dimension, the properties
most often used are size proxies, and the sample sizes used cannot distinguish among 14
collinear predictors.

\subsection{Reconciliation with prior work}

\citet{gutman2013} are not contradicted; they are explained. Their fixed-$n$ design sets
$\CV=0$, makes the bound vacuous, and creates the only conditions under which the indices
can be told apart --- which is why they found real differences. The subsequent literature
restored size variation, adopted additive properties, and reported the resulting size
correlation as predictive success. \citet{randic1994}'s orthogonalization argument stands
and licenses continued interest in intercorrelated descriptors \emph{used in combination};
it does not license single-index comparative rankings. \citet{roy2014}'s warning about
$R^2$ concerns external validation and applies here with additional force, but is
orthogonal to the size-proxy mechanism.

\subsection{A checklist for editors and referees}

The following can be applied without specialist knowledge of chemical graph theory, and
instantiates rather than extends existing QSAR validation practice.
\begin{enumerate}\itemsep2pt
\item \textbf{Report $\Delta R^2$ over a heavy-atom-count null} for every index--property
pair. If $\Delta R^2\approx0$, the result concerns molecular size and should be described
that way.
\item \textbf{Report $\CV(n)$} for the molecule set. Without it, $R^2$ is uninterpretable.
\item \textbf{State whether the property is additive.} Molecular weight, molar
refractivity, molar volume and polarizability are; correlations with them are not evidence
about an index.
\item \textbf{Do not crown a best index at $N<30$} without a resampling analysis of winner
stability.
\item \textbf{Compare $R^2$ against a permutation null computed at the study's own sample
size and fit order.} Indicative 95th percentiles at $N\approx15$: linear $\approx0.35$,
quadratic $\approx0.50$, cubic $\approx0.63$. Raising the fit order raises this bar.
\item \textbf{Report the provenance of property values} --- measured, or estimated by an
additive group-contribution scheme. Regressing an additive estimate on an additive
descriptor is circular.
\end{enumerate}

\subsection{Limitations}

The boiling-point compilation is a handbook aggregation that may include estimated as well
as measured values; the melting-point set is curated experimental data and should be
treated as primary. The CAS-to-structure join through PubChem has an unmeasured mismatch
rate. The maximum-degree $\le4$ filter, which we had expected to bias the sample, excludes
exactly one molecule of $8{,}584$. Our battery covers degree-based indices only ---
distance-based, eccentricity-based and neighbourhood-degree-based families are not tested,
though \eqref{eq:factor} suggests analogous arguments. The drug set of 38 molecules was
assembled by hand and is not a random sample of drug space; all headline claims are
replicated on the larger unbiased sets.

\section*{Declaration of generative AI use}

Portions of this work were carried out with the assistance of Claude, a large language
model developed by Anthropic. Its contribution was substantial and is described here in
full so that readers and referees can weight it appropriately.

The model was used to: survey the prior literature and identify the relevant precedents,
including \citet{gutman2013} and \citet{randic1994}; propose and formalise the profile
decomposition \eqref{eq:factor} and the bound \eqref{eq:bound}--\eqref{eq:spectral};
identify and prove the corollary \eqref{eq:corollary}; write all analysis and figure-generation
code; execute the computational experiments reported in Section~4; and draft substantial
portions of the manuscript text.

The research direction, the choice of which lines of enquiry to pursue and abandon, and
all decisions about scope, framing and claim strength were made by the author. Two of the
model's initial hypotheses --- a worst-case sandwich bound, and a claim that heavy-atom
count would outperform every index out of sample --- were falsified during the work and
retracted; the results reported here reflect the surviving claims. An adversarial review
pass, also model-assisted, identified and corrected a mischaracterised null distribution
and one factual error in an earlier draft.

The author has verified the reported results, takes full responsibility for the content of
this paper, and affirms that the generative AI tool is not and cannot be an author. All
code and data required to reproduce every number and figure are provided as supplementary
material, and independent reproduction is encouraged.

\section*{Reproducibility}

All datasets are public. Analysis code, molecule lists and figure scripts are provided as
supplementary material; the full pipeline runs in a few minutes on a laptop using RDKit
and NumPy.

\appendix
\section{Per-index bound values}

\begin{table}[h]\centering\small
\caption{Spectral bound against observed $r^2(\TI_f,m)$, melting-point set
($N=8{,}583$, $\CV(m)=0.417$). Sorted by slack.}
\begin{tabular}{lrrrr}
\toprule
index & $\tau$ (exact) & $\tau$ (bound) & $1-\tau^2_{\text{spec}}$ & observed $r^2$\\
\midrule
ABC & 0.0605 & 0.0608 & 0.9963 & 0.99663\\
GA & 0.0613 & 0.0614 & 0.9962 & 0.99648\\
sum-connectivity & 0.0924 & 0.0929 & 0.9914 & 0.99231\\
SDD & 0.2680 & 0.2686 & 0.9279 & 0.92885\\
harmonic & 0.1817 & 0.1828 & 0.9666 & 0.96868\\
Randi\'c & 0.1567 & 0.1585 & 0.9749 & 0.97974\\
$M_1$ & 0.1975 & 0.1986 & 0.9606 & 0.97271\\
Sombor & 0.2235 & 0.2242 & 0.9497 & 0.96221\\
ISI & 0.1739 & 0.1761 & 0.9690 & 0.98370\\
AZI & 0.2289 & 0.2316 & 0.9464 & 0.97101\\
$M_2$ & 0.3836 & 0.3867 & 0.8504 & 0.92675\\
HM & 0.4183 & 0.4204 & 0.8232 & 0.90313\\
F & 0.4669 & 0.4683 & 0.7807 & 0.87174\\
ReZG3 & 0.6405 & 0.6447 & 0.5844 & 0.83514\\
\bottomrule
\end{tabular}
\end{table}

\begin{thebibliography}{99}
\bibitem[Delaney(2004)]{delaney2004} Delaney, J.S. (2004) ESOL: estimating aqueous
solubility directly from molecular structure. \emph{J. Chem. Inf. Comput. Sci.}
\textbf{44}(3), 1000--1005. \texttt{doi:10.1021/ci034243x}
\bibitem[Furtula \& Gutman(2015)]{furtula2015} Furtula, B., Gutman, I. (2015) A forgotten
topological index. \emph{J. Math. Chem.} \textbf{53}, 1184--1190.
\bibitem[Gutman(2021)]{gutman2021} Gutman, I. (2021) Geometric approach to degree-based
topological indices: Sombor indices. \emph{MATCH Commun. Math. Comput. Chem.}
\textbf{86}(1), 11--16.
\bibitem[Gutman \& Trinajsti\'c(1972)]{gutman1972} Gutman, I., Trinajsti\'c, N. (1972)
Graph theory and molecular orbitals. Total $\pi$-electron energy of alternant
hydrocarbons. \emph{Chem. Phys. Lett.} \textbf{17}(4), 535--538.
\bibitem[Gutman \& To\v{s}ovi\'c(2013)]{gutman2013} Gutman, I., To\v{s}ovi\'c, J. (2013)
Testing the quality of molecular structure descriptors. Vertex-degree-based topological
indices. \emph{J. Serb. Chem. Soc.} \textbf{78}(6), 805--810.
\texttt{doi:10.2298/JSC121002134G}
\bibitem[Randi\'c(1975)]{randic1975} Randi\'c, M. (1975) On characterization of molecular
branching. \emph{J. Am. Chem. Soc.} \textbf{97}(23), 6609--6615.
\textbf{[VERIFY TITLE: cited in the literature both as ``On characterization of molecular
branching'' and as ``Characterization of molecular branching''; check the ACS record.]}
\bibitem[Randi\'c et al.(1994)]{randic1994} Randi\'c, M., et al. (1994) Graphical bond
orders: novel structural descriptors. \emph{J. Chem. Inf. Comput. Sci.} \textbf{34}(2).
\textbf{[VERIFY: full author list and page range unconfirmed; article confirmed to span
p.~405. Check the ACS record.]}
\bibitem[Roy \& Gramatica(2014)]{roy2014} Roy, K., Gramatica, P. (2014) Beware of $R^2$:
simple, unambiguous assessment of the prediction accuracy of QSAR and QSPR models.
\emph{J. Chem. Inf. Model.} \textbf{54}(7), 1696--1701.
\end{thebibliography}

\end{document}
